\providecommand{\GordanProjectName}{Paul Gordan, De linea geodetica (1862)}
\newcommand{\gq}[1]{\glqq#1\grqq}
\newcommand{\tg}{\operatorname{tg}}
\newcommand{\arc}{\operatorname{arc}}
\newcommand{\dd}{\,d}
\newcommand{\GordanDitto}{\mbox{\normalfont\ttfamily\char34\kern0.45em\char34}}
\newcommand{\caput}[1]{\par\vspace{0.35em}\begin{center}{\large\bfseries CAPUT #1.}\end{center}\vspace{-0.75em}}
\newcommand{\sect}[1]{\par\vspace{0.28em}\begin{center}{\bfseries\S\ #1.}\end{center}\vspace{-0.65em}}
\newcommand{\GordanUnitStart}[2]{%
  \clearpage\thispagestyle{empty}%
  \phantomsection\hypertarget{PU#1}{}%
  \pdfbookmark[1]{Physical unit #1 -- #2}{bookmark-PU#1}%
}
\newcommand{\GordanTopFolio}[1]{%
  \begin{center}\textbf{#1}\\[-0.45em]\rule{3.0cm}{0.35pt}\end{center}\vspace{0.1em}%
}
\newcommand{\GordanBottomFolio}[1]{\par\vfill\hfill\textbf{#1}\par}
\newcommand{\GordanOrnament}{\begin{center}\rule{2.2cm}{0.35pt}\kern-1.25cm\raisebox{-0.12ex}{\scriptsize$\blacklozenge$}\end{center}}
